% This is text associated with the open textbook "Open Electromagnetics"
% See immediately below for author, date
% This work is licensed under the Creative Commons Attribution-ShareAlike 4.0 International License. (CC BY-SA 4.0) 
% To view a copy of the license, visit http://creativecommons.org/licenses/by-sa/4.0/

%ORB TITLE Power Radiated by an Electrically-Short Dipole
%ORB AUTHOR 0        # S.W. Ellingson, Virginia Tech (ellingson@vt.edu)
%ORB DATE 20191126   # yyyymmdd
%ORB ABSTRACT 

\label{m0207_Total_Power_Radiated_by_an_Electrically-Short_Dipole} % <-- Must match name of this file and module directory name.  Don't mess with this!
\index{antenna!electrically-short dipole (ESD)|(}

In this section, we determine the total power radiated by an electrically-short dipole (ESD) antenna in response to a sinusoidally-varying current applied to the antenna terminals.  This result is both useful on its own and necessary as an intermediate result in determining the impedance of the ESD.
The ESD is introduced in Section~\ref{m0198_Radiation_from_an_Electrically-Short_Dipole}, and a review of that section is suggested before attempting this section.

In Section~\ref{m0198_Radiation_from_an_Electrically-Short_Dipole}, it is shown that the electric field intensity in the far field of a $\hat{\bf z}$-oriented ESD located at the origin is 
%
\begin{equation}
\widetilde{\bf E}({\bf r}) \approx  \hat{\bf \theta} j \eta \frac{I_0\cdot\beta L}{8\pi} ~\left(\sin\theta\right) ~\frac{e^{-j\beta r}}{r}
\label{m0207_eE}
\end{equation}
%
where 
${\bf r}$ is the field point,
$I_0$ is a complex number representing the peak magnitude and phase of the sinusoidally-varying terminal current,
$L$ is the length of the ESD,
and $\beta$ is the phase propagation constant $2\pi/\lambda$ where $\lambda$ is wavelength.
Note that $L\ll\lambda$ since this is an ESD. 
Also note that Equation~\ref{m0207_eE} is valid only for the far-field conditions $r\gg L$ and $r\gg \lambda$, and presumes propagation in simple (linear, homogeneous, time-invariant, isotropic) media with negligible loss.

Given that we have already limited scope to the far field, it is reasonable to approximate the electromagnetic field at each field point ${\bf r}$ as a plane wave propagating radially away from the antenna; i.e., in the $\hat{\bf r}$ direction.  Under this assumption, the time-average power density is 
%
\begin{equation}
{\bf S}({\bf r}) = \hat{\bf r}\frac{\left|\widetilde{\bf E}({\bf r})\right|^2}{2\eta} 
\end{equation}
%
where $\eta$ is the wave impedance of the medium.
The total power $P_{rad}$ radiated by the antenna is simply ${\bf S}({\bf r})$ integrated over any closed surface $\mathcal{S}$ that encloses the antenna.
Thus:
%
\begin{equation}
P_{rad} = \oint_{\mathcal{S}} { {\bf S}({\bf r}) \cdot d{\bf s} }
\end{equation}
%
where $d{\bf s}$ is the outward-facing differential element of surface area.
In other words, power density (W/m$^2$) integrated over an area (m$^2$) gives power (W).
Anticipating that this problem will be addressed in spherical coordinates, we note that
%
\begin{equation}
d{\bf s} = \hat{\bf r} r^2~\sin\theta~d\theta~d\phi
\end{equation}
% 
and subsequently:
%
\begin{flalign}
P_{rad} &= \int_{\theta=0}^{\pi} \int_{\phi=0}^{2\pi} { \left(\hat{\bf r}\frac{\left|\widetilde{\bf E}({\bf r})\right|^2}{2\eta} \right) \cdot \left(\hat{\bf r} r^2~\sin\theta~d\theta~d\phi\right) } \nonumber \\
        &= \frac{1}{2\eta}\int_{\theta=0}^{\pi} \int_{\phi=0}^{2\pi} { \left|\widetilde{\bf E}({\bf r})\right|^2 r^2 \sin\theta~d\theta~d\phi } \label{m0207_ePrad} 
\end{flalign}
%
Returning to Equation~\ref{m0207_eE}, we note:
%
\begin{equation}
\left|\widetilde{\bf E}({\bf r})\right|^2 \approx \eta^2 \frac{\left|I_0\right|^2\left(\beta L\right)^2}{64\pi^2} ~\left(\sin\theta\right)^2 ~\frac{1}{r^2}
\label{m0207_eE2}
\end{equation}
%
Substitution into Equation~\ref{m0207_ePrad} yields:
%
\begin{equation}
P_{rad} \approx \eta \frac{\left|I_0\right|^2\left(\beta L\right)^2}{128\pi^2}\int_{\theta=0}^{\pi} \int_{\phi=0}^{2\pi} { \sin^3 \theta~d\theta~d\phi } 
\end{equation}
%
Note that this may be factored into separate integrals over $\theta$ and $\phi$.  The integral over $\phi$ is simply $2\pi$, leaving:
%
\begin{equation}
P_{rad} \approx \eta \frac{\left|I_0\right|^2\left(\beta L\right)^2}{64\pi}\int_{\theta=0}^{\pi} { \sin^3 \theta~d\theta } 
\end{equation}
%
The remaining integral is equal to $4/3$, leaving:
%
\begin{equation}
\boxed{
P_{rad} \approx \eta \frac{\left|I_0\right|^2\left(\beta L\right)^2}{48\pi}
}
\label{m0207_ePT}
\end{equation}
%
This completes the derivation, but it is useful to check units.  
Recall that $\beta$ has SI base units of rad/m, so $\beta L$ has units of radians.
This leaves $\eta\left|I_0\right|^2$, which has SI base units of $\Omega$ $\cdot$ A$^2$ $=$ W, as expected.

%%%%%%%%%%%%%%%%%%%%%%%%%%%%%%%%%%%%%%%%%%%%%%%
%ORB HIGHLIGHT
\begin{mdframed}[backgroundcolor=yellow!20,nobreak=true] 
The power radiated by an ESD in response to the current $I_0$ applied at the terminals is given by Equation~\ref{m0207_ePT}.
\end{mdframed}
%%%%%%%%%%%%%%%%%%%%%%%%%%%%%%%%%%%%%%%%%%%%%%%

Finally, it is useful to consider how various parameters affect the radiated power.
First, note that the radiated power is proportional to the square of the terminal current.
Second, note that the product $\beta L=2\pi L/\lambda$ is the electrical length
\index{electrical length}
of the antenna; that is,  the length of the antenna expressed in radians, where $2\pi$ radians is one wavelength.
Thus, we see that the power radiated by the antenna increases as the square of electrical length.

%%%%%%%%%%%%%%%%%%%%%%%%%%%%%%%%%%%%%%%%%%%%%%%
%ORB EXAMPLE
~\\ \begin{example} 
\begin{mdframed}[backgroundcolor=brown!20,linecolor=brown!20]\raggedright
Power radiated by an ESD. 
\\~\\
A dipole is 10~cm in length and is surrounded by free space.
A sinusoidal current having frequency 30~MHz and peak magnitude 100~mA is applied to the antenna terminals.
What is the power radiated by the antenna?

{\bf Solution.}
If no power is dissipated within the antenna, then all power is radiated.
The wavelength $\lambda=c/f \cong 10$~m, so $L \cong 0.01\lambda$.
This certainly qualifies as electrically-short, so we may use Equation~\ref{m0207_ePT}.
In the present problem, $\eta\cong 376.7~\Omega$ (the free space wave impedance), $I_0=100$~mA, and $\beta=2\pi/\lambda\cong 0.628$~rad/m.
Thus, we find that the radiated power is $\approx$ \underline{98.6~$\mu$W}.
\end{mdframed}
% note figures will need to go between "\end{mdframed}" and "\end{example}"
\end{example}
%%%%%%%%%%%%%%%%%%%%%%%%%%%%%%%%%%%%%%%%%%%%%%%

\index{antenna!electrically-short dipole (ESD)|)}

